\documentclass[11pt,a4paper]{article}
\usepackage[italian]{babel}
\usepackage[margin=2.5cm]{geometry}
\usepackage{parskip}
\usepackage{amsmath, amssymb, amsthm}
\usepackage[most]{tcolorbox}
\usepackage{array}
\usepackage{tikz}
\usetikzlibrary{decorations.pathreplacing}
\usepackage{hyperref}

% --- Riquadri con bordo nero, senza numero ---
% Uso: \begin{definizione} ... \end{definizione}
%      \begin{teorema}[Nome] ... \end{teorema}  (il nome è facoltativo)
\tcbset{nota/.style={colback=white, colframe=black, boxrule=0.5pt,
  sharp corners}}
\NewTColorBox{definizione}{}{nota,
  before upper={\textbf{Definizione.}\ }}
\NewTColorBox{teorema}{o}{nota,
  before upper={\textbf{Teorema\IfValueT{#1}{ (#1)}.}\ }}
\NewTColorBox{proposizione}{o}{nota,
  before upper={\textbf{Proposizione\IfValueT{#1}{ (#1)}.}\ }}
\NewTColorBox{assioma}{o}{nota, boxrule=1pt,
  before upper={\textbf{Assioma\IfValueT{#1}{ (#1)}.}\ }}

% --- Ambienti semplici (senza riquadro) ---
\theoremstyle{definition}
\newtheorem*{esempio}{Esempio}
\newtheorem*{osservazione}{Osservazione}
\newtheorem*{esercizio}{Esercizio}

% V e F nelle tavole di verità
\newcommand{\V}{\textsf{V}}
\newcommand{\F}{\textsf{F}}

\title{Esempi sugli estremi e insiemi non limitati}
\author{Marco Cerbella}
\date{Lezione 3 -- 1 ottobre 2026}

\begin{document}
\maketitle

\section{Un altro esempio: $B = \left\{ \dfrac{n}{n+1} \mid n \in \mathbb{N} \right\}$}

\begin{esempio}
Sia $B = \{x_n = \frac{n}{n+1},\ n \in \mathbb{N}\} = \{0, \frac12, \frac23, \frac34, \dots\}$. Poiché $x_n \leq 1$ per ogni $n$ (infatti $n \leq n + 1$) e $x_n \geq 0$, si ha $B \subseteq [0, 1]$: $B$ è limitato. In particolare, essendo limitato superiormente, per (R4) esiste $\sup B$.

\textbf{Minimo.} $x_0 = 0 \in B$ e $x_n \geq 0$ per ogni $x_n \in B$: quindi $0 = \min B$.

\textbf{Estremo superiore.} Dimostriamo che $\sup B = 1$, cioè che:
\begin{enumerate}
    \item $x_n \leq 1$ per ogni $x_n \in B$ (già visto);
    \item per ogni $\varepsilon > 0$ esiste $x_n \in B$ tale che $1 - \varepsilon < x_n$.
\end{enumerate}
Per assurdo neghiamo la 2: esiste $\bar\varepsilon > 0$ tale che $1 - \bar\varepsilon \geq x_n = \frac{n}{n+1}$ per ogni $n \in \mathbb{N}$. Allora
\[
(n+1)(1 - \bar\varepsilon) \geq n \iff n + 1 - \bar\varepsilon n - \bar\varepsilon \geq n
\iff 1 - \bar\varepsilon \geq \bar\varepsilon n
\iff n \leq \frac{1 - \bar\varepsilon}{\bar\varepsilon} \quad \forall\, n \in \mathbb{N}.
\]
Questo è assurdo, perché $n$ può essere grande a piacere. Quindi $1 = \sup B$.

Poiché $1 \notin B$ (si ha $x_n < 1$ per ogni $n$), $B$ non ha massimo.
\end{esempio}

\section{Insiemi non limitati}

Un insieme $A$ non sempre ammette maggioranti o minoranti.

\begin{definizione}
Si pone
\begin{itemize}
    \item $\sup A = +\infty$ se $A$ non è limitato superiormente, cioè
          $\forall\, L \ \exists\, x \in A : x > L$;
    \item $\inf A = -\infty$ se $A$ non è limitato inferiormente, cioè
          $\forall\, \ell \ \exists\, x \in A : x < \ell$.
\end{itemize}
\end{definizione}

\begin{esempio}
\leavevmode
\begin{itemize}
    \item $A = [0, +\infty)$ è limitato inferiormente ma non superiormente;
    \item $A = (-\infty, 0]$ è limitato superiormente ma non inferiormente;
    \item $[-2, +\infty)$ è limitato inferiormente ma non superiormente;
    \item $(-\infty, 5]$ è limitato superiormente ma non inferiormente;
    \item $\,]2, 1000[$ è limitato: $2 = \inf\,]2, 1000[$ e
          $1000 = \sup\,]2, 1000[$, ma l'insieme non ha né minimo né massimo
          (perché $2, 1000 \notin \,]2, 1000[$);
    \item $\mathbb{N}$ è limitato inferiormente ma non superiormente.
\end{itemize}
\end{esempio}

\begin{esercizio}
Quale tra i seguenti insiemi ha sia massimo sia minimo?
(1) $]2, 9[$ \qquad (2) $[0, 5[$ \qquad (3) $[0, 1[ \,\cup \{7\}$ \qquad (4) $]0, 2] \cup \{5\}$
\textbf{Risposta: 3.} $[0, 1[ \,\cup \{7\}$ ha minimo $0$ e massimo $7$. Gli altri: $]2, 9[$ non ha né minimo né massimo; $[0, 5[$ ha minimo $0$ ma non massimo; $]0, 2] \cup \{5\}$ ha massimo $5$ ma non minimo.
\end{esercizio}

\end{document}
